% answers_short.tex —— 1977 年·河南卷 · 简短答案速查
% 由 tools/make_short_answers.py 自 answers.tex 的 \gksolhead/\gkanswer 提取，
% 并与 1977/1977_数学真题_河南_答案要点.md 逐题比对后生成，勿手改（改源 MD 后重跑）。

\begin{gknotice}
\textbf{使用说明}：本节只给最终答案便于快速核对，完整推导见解析版。答案为编写者独立推导，\textbf{非官方标准答案}。
\end{gknotice}

\gkpart{题目}
\begin{itemize}[leftmargin=2.6em]
  \item[\textbf{1.}] $0$。
  \item[\textbf{2.}] $1$。
  \item[\textbf{3.}] $x=3$ 或 $x=-1$。
  \item[\textbf{4.}] 见解析版
  \item[\textbf{5.}] $BG\perp DF$。
  \item[\textbf{6.}] 见解析版
  \item[\textbf{7.}] $V=\dfrac12L^3\cos^2\alpha\sin\alpha\sin\beta$。
  \item[\textbf{8.}] (1) $f_1(x)=-\dfrac14x^2+\dfrac12x+\dfrac{11}{4}=-\dfrac14(x-1)^2+3$；$f_2(x)=\dfrac14x^2-\dfrac12x+\dfrac54=\dfrac14(x-1)^2+1$；(2) $f_1(x)$ 在 $x=1$ 处取极大值 $3$；$f_2(x)$ 在 $x=1$ 处取极小值 $1$；(3) 见解析版
  \item[\textbf{9.}] (1) $\begin{aligned}M&=\left(\dfrac{-3\sqrt2\,k+3\sqrt{k^2-1}}{1+k^2},\ \dfrac{3\sqrt2+3k\sqrt{k^2-1}}{1+k^2}\right),\\ N&=\left(\dfrac{-3\sqrt2\,k-3\sqrt{k^2-1}}{1+k^2},\ \dfrac{3\sqrt2-3k\sqrt{k^2-1}}{1+k^2}\right).\end{aligned}$；(2) 当 $M\text{、}N$ 重合时 $k=\pm1$；此时直线与圆相切，这样的直线有两条：$y=x+3\sqrt2$ 与 $y=-x+3\sqrt2$，其夹角为 $90^\circ$
  \item[\textbf{10.}] (1) 第一根电线杆距电管所 $1000$ 米；(2) 共栽 $30$ 根
\end{itemize}
